% Source: https://artofproblemsolving.com/community/c3h2211776_problem
% Stable identity: aops-topic-2211776-q01
\textbf{(a)} By the AM-GM inequality for the $n$ positive numbers $1,2,\dots,n$,
\[\dfrac{1+2+\cdots+n}{n}\ge\sqrt[n]{1\cdot2\cdots n}.\]
The left side is $\dfrac{n(n+1)}{2n}=\dfrac{n+1}{2}$, so raising to the $n$th power gives $\left(\dfrac{n+1}{2}\right)^n\ge n!$ (with equality only for $n=1$). $\blacksquare$

\textbf{(b)} The $n$th root of $\left(\dfrac{n+1}{2}\right)^n$ is $\dfrac{n+1}{2}$, but the $n$th root of $n!$ is $\sqrt[n]{n!}$, not $(n-1)!$; the student treated taking a root as if it removed one factor. The correct consequence is only $\dfrac{n+1}{2}\ge\sqrt[n]{n!}$.

The conclusion is false: for $x=3$ it claims $\dfrac{5}{2}\ge3!=6$. In fact $\dfrac{x+2}{2}\ge x!$ holds only for $x=1$ and $x=2$ (for $x=2$, $2\ge2$), and for $x\ge3$ we have $x!\ge2x>\dfrac{x+2}{2}$.
